\chapter{Panel Data}\label{ch:b3:panel-data}

Richer countries emit more carbon dioxide per person. Across 200 countries,
a country whose income per person is one log point higher emits, on average,
2.7 more tonnes per person. % figdata: co2_stats (between 2.71597)
Follow each country over time instead, and the slope shrinks to about one
tonne. % figdata: co2_stats (within 0.999459)
Both numbers come from the same 6{,}419 observations. The first compares
countries with each other; the second compares each country with itself in
other years. They differ because countries differ in many fixed ways --- coal
under the ground, a cold winter, a car-built city --- that go with both their
income and their emissions. A data set that follows the same units over time
can take those fixed differences out, without measuring them. This chapter is
about how, at what price, and what the method cannot remove.

\begin{openingchart}{Carbon dioxide emissions per person against income per
person, 1990--2022, 200 countries. Left: across countries, each dot a
country's average over the years (GDP per person in constant 2015 dollars, log
scale). Right: inside countries, each dot the mean of one twentieth of the
6{,}419 country-years, after removing each country's own averages. Lines:
least-squares fits. Palau, at 104 tonnes, lies above the left panel's scale.}% figdata: between (PLW 104.489)
{World Bank, World Development Indicators
(EN.GHG.CO2.PC.CE.AR5, tonnes of CO$_2$ per person; NY.GDP.PCAP.KD); our
computation}
\begin{tikzpicture}
\begin{groupplot}[group style={group size=2 by 1, horizontal sep=14mm},
  omscatter, height=0.5\linewidth]
\nextgroupplot[xmode=log, width=0.5\linewidth, xmin=100, xmax=200000,
  xtick={100,1000,10000,100000}, log ticks with fixed point, ymin=0, ymax=50,
  xlabel={GDP per person (2015 \$, log scale)}, ylabel={CO$_2$ per person (tonnes)},
  title={\small Between countries}]
  \addplot table[x=gdp, y=co2] {figdata/bachelor-3/02-panel-data/out/between.csv};
  \addplot[omThm, thick, no marks, sharp plot] table[x=gdp, y=co2] {figdata/bachelor-3/02-panel-data/out/between_line.csv};
\nextgroupplot[width=0.5\linewidth, xmin=-0.8, xmax=0.8, ymin=-1.2, ymax=1.2,
  xlabel={Log income, minus the country's mean}, ylabel={CO$_2$, minus the country's mean},
  title={\small Within countries},
  x tick label style={/pgf/number format/.cd, fixed, precision=1},
  y tick label style={/pgf/number format/.cd, fixed, precision=1}]
  \addplot table[x=dly, y=dco2] {figdata/bachelor-3/02-panel-data/out/within_bins.csv};
  \addplot[omThm, thick, no marks, sharp plot] table[x=dly, y=dco2] {figdata/bachelor-3/02-panel-data/out/within_line.csv};
\end{groupplot}
\end{tikzpicture}
\end{openingchart}

The left panel is the familiar cross-section. The rich oil producers sit high
above the line, and several rich countries with nuclear or hydroelectric power
sit below it. The right panel removes every country's average income and
average emissions, so that a country only ever meets itself: a year in which
it was ten per cent richer than its own average is compared with its own
average emissions. That slope, 1.0, is less than half of the slope across
countries, 2.7. % figdata: co2_stats
Note the horizontal scale of the right panel. Almost all the variation in
income is between countries. What is left inside them is a band of a few
tenths of a log point, and the within slope is estimated from that narrow
band alone.

\begin{recall}
The Year 1 volume introduced panel data as a data shape: the same units observed
at several dates. The Year 2 volume proved the Frisch--Waugh--Lovell theorem (the
coefficient on $x$ in a regression on $x$ and $W$ is the slope of $y$ on the
part of $x$ that $W$ does not predict), derived omitted-variable bias, and
defined cluster-robust standard errors and the Moulton factor.
\Cref{ch:b3:instrumental-variables} met an \omterm{def:b3:instrumental-variables:endogenous}{endogenous regressor} and one cure,
the instrument.
\end{recall}

\begin{notation}[Symbols of this chapter]\label{not:b3:panel-data:symbols}
A unit (a country, a firm) is $i = 1, \dots, n$ and a period $t = 1, \dots, T$;
here $T$ is the length of the panel, not taxes. The unit effect is $a_i$ and a
period effect $\mu_t$. A bar over a variable is its mean over the periods of
one unit, $\bar x_i = \frac1T \sum_t x_{it}$, and a double dot the deviation
from it, $\ddot x_{it} = x_{it} - \bar x_i$.
\end{notation}

\section{The same units over time}\label{sec:b3:panel-data:units}

The model of this chapter is
\begin{equation}\label{eq:b3:pd:model}
  y_{it} = \beta x_{it} + a_i + u_{it}, \qquad \E[u_{it} \mid x_{i1}, \dots, x_{iT}, a_i] = 0.
\end{equation}
The unit effect $a_i$ collects everything about unit $i$ that does not change
over the panel and affects $y$. It is not observed. It may be correlated with
$x$, and in most applications it is: the reason to use a panel is that it is.
The condition on $u_{it}$, called strict exogeneity, says that the remaining
error is unrelated to the regressor in every period, past and future.

\begin{definition}[Pooled regression]\label{def:b3:panel-data:pooled}
\index{pooled regression}
The \emph{pooled regression} is the OLS regression of $y_{it}$ on $x_{it}$ (and a
constant) over all $nT$ observations, ignoring which unit each comes from.
\end{definition}

The \omterm{def:b3:panel-data:pooled}{pooled regression} leaves $a_i$ in its error. By the omitted-variable formula,
its slope converges to $\beta + \Cov(x_{it}, a_i)/\Var(x_{it})$. The figure
below shows the pattern with three units: within each unit the relation has
slope $\beta$, but the units with high $x$ also have high $a_i$, and the pooled
line, steeper than any of them, joins the clouds.

\begin{omfigure}
\begin{tikzpicture}[econ, scale=0.85]
  \draw[econ axes] (0,5.4) node[above] {$y$} |- (8,0) node[right] {$x$};
  \foreach \cx/\cy/\col in {1.6/0.9/omDef, 3.9/2.7/omThm, 6.2/4.5/omMeth} {
    \draw[\col, very thick] (\cx-0.9,\cy-0.225) -- (\cx+0.9,\cy+0.225);
    \foreach \dx/\dy in {-0.7/-0.12, -0.3/0.14, 0.1/-0.1, 0.5/0.12, 0.8/-0.08} {
      \fill[\col] (\cx+\dx,\cy+0.25*\dx+\dy) circle (1.6pt);
    }
  }
  \draw[curve, dashed] (0.6,0.12) -- (7.3,5.35) node[right] {pooled};
  \node[font=\footnotesize, omDef, anchor=north] at (1.6,0.45) {unit 1};
  \node[font=\footnotesize, omThm, anchor=north] at (3.9,2.25) {unit 2};
  \node[font=\footnotesize, omMeth, anchor=north] at (6.2,4.05) {unit 3};
\end{tikzpicture}
\par\smallskip{\small Three units with the same slope and different
intercepts $a_i$. The units with larger $x$ have larger $a_i$, so the pooled
line (dashed) is steeper than the slope inside every unit.}
\end{omfigure}

\section{Fixed effects}\label{sec:b3:panel-data:fixed-effects}

\begin{definition}[Unit fixed effect, within estimator]\label{def:b3:panel-data:within}
\index{unit fixed effect}\index{within estimator}
In model~\eqref{eq:b3:pd:model}, treating $a_i$ as a parameter to be estimated
for each unit makes it a \emph{unit fixed effect}. The \emph{within estimator}
(or fixed-effects estimator) of $\beta$ is the OLS slope of the deviations
$\ddot y_{it}$ on the deviations $\ddot x_{it}$:
\[
  \hat\beta_{\mathrm{FE}} = \frac{\sum_i \sum_t \ddot x_{it}\, \ddot y_{it}}{\sum_i \sum_t \ddot x_{it}^2}.
\]
\end{definition}

Averaging model~\eqref{eq:b3:pd:model} over $t$ gives $\bar y_i = \beta \bar
x_i + a_i + \bar u_i$; subtracting it from the model removes $a_i$:
$\ddot y_{it} = \beta \ddot x_{it} + \ddot u_{it}$. Under strict exogeneity,
$\ddot u_{it}$, which contains $\bar u_i$ and hence the errors of every period,
is uncorrelated with $\ddot x_{it}$, and the \omterm{def:b3:panel-data:within}{within estimator} is consistent as
$n$ grows. Whatever $a_i$ is, and however it is related to $x$, it is gone.

\begin{proposition}[Within estimator and unit dummies]\label{prop:b3:panel-data:dummies}
The \omterm{def:b3:panel-data:within}{within estimator} equals the coefficient on $x_{it}$ in the OLS regression of
$y_{it}$ on $x_{it}$ and $n$ unit dummies $d^{(1)}_{it}, \dots, d^{(n)}_{it}$,
where $d^{(j)}_{it} = 1$ if $i = j$.
\end{proposition}

\begin{proof}
By the Frisch--Waugh--Lovell theorem, the coefficient on $x$ is the slope of $y$
on the residual of $x$ after a regression of $x$ on the dummies. Regressing on the
dummies fits one constant per unit, and the least-squares constant for unit $i$
is the unit's mean, $\bar x_i$. The residual is $x_{it} - \bar x_i = \ddot
x_{it}$. The theorem also allows $y$ to be replaced by its own residual, $\ddot
y_{it}$, so the coefficient is $\sum \ddot x_{it} \ddot y_{it} / \sum \ddot
x_{it}^2$.
\end{proof}

The dummies give the same slope with $n$ more parameters; the within transform
gives it without them. Either way, the degrees of freedom of the residual
variance are $nT - n - 1$, a point that software using the transform must
handle.

\begin{definition}[First-difference estimator]\label{def:b3:panel-data:first-difference}
\index{first-difference estimator}
The \emph{first-difference estimator} of $\beta$ is the OLS slope of $\Delta
y_{it} = y_{it} - y_{i,t-1}$ on $\Delta x_{it} = x_{it} - x_{i,t-1}$, for $t =
2, \dots, T$.
\end{definition}

Differencing also removes $a_i$. With two periods the two estimators are
identical (\cref{exo:b3:panel-data:10}). With more, they weight the data
differently and are both consistent under strict exogeneity; they differ when
the errors are serially correlated, when $x$ is measured with error, and when
the model is wrong. A gap between them is a warning worth reading.

\begin{definition}[Two-way fixed effects]\label{def:b3:panel-data:twfe}
\index{two-way fixed effects}
The \emph{two-way fixed effects} model adds a period effect to
model~\eqref{eq:b3:pd:model}: $y_{it} = \beta x_{it} + a_i + \mu_t + u_{it}$.
It is estimated with unit and period dummies, or, in a balanced panel, by the
slope of $y_{it} - \bar y_i - \bar y_t + \bar y$ on the same transform of $x$.
\end{definition}

\Cref{lst:b3:panel-data:co2} computes the slopes behind the opening chart.

\omcode{code/bachelor-3/02-panel-data/listings/pooled_vs_within.py}{1}{24}{Pooled, between, within, dummies and two-way slopes, on the World Development Indicators.}\label{lst:b3:panel-data:co2}

\omoutput{code/bachelor-3/02-panel-data/outputs/pooled_vs_within.txt}{Tonnes of CO$_2$ per person per log point of income.}\label{out:b3:panel-data:co2}

The pooled and between slopes are close, 2.756 and 2.716, because most of the
variation is between countries; only 3.6~\% of the variance of log income is
within. % output: pooled_vs_within
The within slope and the slope with 200 country dummies agree to every printed
digit, 0.999, as \cref{prop:b3:panel-data:dummies} requires. Adding year
effects doubles it, to 2.068. % output: pooled_vs_within
The year effects take out what all countries share in a year: world income
rose over the period while some large rich emitters cut their emissions per
person, so a common time path pulled the one-way within slope down. Which
slope answers which question is the subject of
\cref{sec:b3:panel-data:identifying}.

\begin{model}{The fixed-effects model}
\assumes Each unit has a constant unobserved level $a_i$, possibly correlated
with $x$; the slope $\beta$ is the same for all units; the remaining error is
unrelated to the regressor in every period (strict exogeneity).
\predicts The within and \omterm{def:b3:panel-data:first-difference}{first-difference estimators} agree, up to sampling
error; adding variables that are fixed for each unit changes nothing;
regressing $y$ on future values of $x$ gives nothing once current $x$ is in.
\tested Not once for all: its predictions are checked in each application,
by comparing within, first-difference and two-way estimates and by the lead
test of \cref{exo:b3:panel-data:12}. One of its working assumptions has been
tested on real data and failed: errors independent over time. In state panels
of female wages, placebo laws came out significant far more often than 5~\%
(\cref{sec:b3:panel-data:clustering}; Bertrand, Duflo and Mullainathan 2004). % ledger: b3pd:bdm04.placebo
\breaks When the confounder changes over time (a new manager, a policy that
moves with the business cycle); when today's $x$ responds to yesterday's $u$
(feedback, as in a dynamic panel, \cref{sec:b3:panel-data:dynamic}); when
the slope differs across units, in which case the \omterm{def:b3:panel-data:within}{within estimator} is a
weighted average of unit slopes, weighted by each unit's variance of $x$.
\end{model}

\section{Random effects}\label{sec:b3:panel-data:random-effects}

The fixed-effects estimator throws away the between variation, which is most of
the data. If $a_i$ were known to be uncorrelated with $x$, nothing would need
throwing away: the \omterm{def:b3:panel-data:pooled}{pooled regression} would be consistent, but its errors $a_i +
u_{it}$ would be correlated within each unit, and a better estimator would use
that structure.

\begin{definition}[Random-effects estimator]\label{def:b3:panel-data:random-effects}
\index{random-effects estimator}
In model~\eqref{eq:b3:pd:model}, suppose $a_i$ has variance $\sigma_a^2$, $u_{it}$
has variance $\sigma_u^2$, and both are uncorrelated with each other and with
$x$. The \emph{random-effects estimator} is the OLS slope after quasi-demeaning,
$y_{it} - \theta \bar y_i$ on $x_{it} - \theta \bar x_i$ (with a constant), where
$\theta = 1 - \sqrt{\sigma_u^2/(\sigma_u^2 + T\sigma_a^2)}$, the variances being
replaced by estimates.
\end{definition}

The weight $\theta$ runs from 0, when there is no unit effect and random effects
is the \omterm{def:b3:panel-data:pooled}{pooled regression}, to 1, when the unit effect dominates and random
effects is the \omterm{def:b3:panel-data:within}{within estimator}. It is the generalised least squares estimator
for this error structure.
\admitted
The proof that quasi-demeaning makes the composite errors uncorrelated with
equal variances is a computation with the $T \times T$ covariance matrix
$\sigma_u^2 I + \sigma_a^2 \mathbf 1 \mathbf 1'$, found in graduate texts on panel
econometrics.

\begin{proposition}[When random effects is consistent]\label{prop:b3:panel-data:re}
Under strict exogeneity of $u$, and with $\theta < 1$, the \omterm{def:b3:panel-data:random-effects}{random-effects
estimator} is consistent if $\Cov(a_i, x_{it}) = 0$ and in general inconsistent
otherwise.
\end{proposition}

\begin{proof}
After quasi-demeaning, the error of the transformed model is $(1 - \theta)a_i +
u_{it} - \theta \bar u_i$. The terms in $u$ are uncorrelated with the transformed
regressor $x_{it} - \theta \bar x_i$ by strict exogeneity. The term $(1 -
\theta)a_i$ is uncorrelated with it if and only if $\Cov(a_i, x_{it} - \theta \bar
x_i) = 0$, which holds when $a_i$ is uncorrelated with $x$ in every period. When
it is not, the transformed error is correlated with the transformed regressor
and the omitted-variable formula gives a bias proportional to $1 - \theta$.
\end{proof}

\omcode{code/bachelor-3/02-panel-data/listings/random_effects.py}{1}{33}{Random against fixed effects, with and without a correlated unit effect.}\label{lst:b3:panel-data:re}

\omoutput{code/bachelor-3/02-panel-data/outputs/random_effects.txt}{True slope 1 in both panels.}\label{out:b3:panel-data:re}

With the unit effect unrelated to $x$, both estimators find the truth, 0.996
and 0.999. When the correlation is 0.45, the fixed-effects estimate stays at
1.003 and random effects moves to 1.149. % output: random_effects
The comparison is the logic of the Hausman test: under the null that $a_i$ is
unrelated to $x$ the two estimates converge to the same value, and a large
gap rejects random effects. In observational economics the null rarely
survives. A useful compromise, due to Mundlak (1978), adds the unit means
$\bar x_i$ to a random-effects regression; the coefficient on $x_{it}$ is then
exactly the within estimate, and the coefficient on $\bar x_i$ measures how
far between and within disagree.

\section{What variation identifies the coefficient}\label{sec:b3:panel-data:identifying}

\begin{definition}[Identifying variation]\label{def:b3:panel-data:identifying-variation}
\index{identifying variation}
The \emph{identifying variation} of an estimator is the part of the variation
of the regressor that the estimator uses: the variation left in $x$ after the
fixed effects and controls are removed, which by the Frisch--Waugh--Lovell
theorem is all the slope depends on.
\end{definition}

The definition turns a statistical fact into a question about the world. The
within slope of the opening chart uses only the 3.6~\% of income variation that
happens inside countries over three decades: % output: pooled_vs_within
booms, recessions, and three decades of growth. Its comparison is between a
country in a good year and the same country in a bad one, and a bad year's
emissions fall for reasons that have to do with the energy used in a recession,
not with the slow change in technology that the between slope reflects. The two
slopes answer different questions, and neither is the long-run effect of
becoming richer on emissions.

\begin{measure}[What a fixed effect absorbs]
A \omterm{def:b3:panel-data:within}{unit fixed effect} absorbs every variable that is constant within the unit
over the panel: geography, climate, legal origin, the energy endowment, a
firm's founding year. Their coefficients cannot be estimated with unit effects,
and a regression that reports one is reporting noise or a collinearity the
software resolved silently. A period effect absorbs everything common to all
units in a period: world prices, global recessions, a change in a survey's
method. What is left --- the \omterm{def:b3:panel-data:identifying-variation}{identifying variation} --- is often small: 3.6~\%
of the variance of log income in the opening chart. % output: pooled_vs_within
Two consequences follow. Measurement error, which is mostly within-unit noise,
becomes a larger share of what is left, so attenuation bias is worse in the
within estimate than in the pooled one. And the standard errors of a
fixed-effects estimate are larger, honestly, because there is less information.
\end{measure}

\section{Errors that move together: clustering in panels}\label{sec:b3:panel-data:clustering}

In a panel, a unit's errors are rarely independent from one year to the next. A
firm that has a good year tends to have a good next year, and a state's wage
level drifts slowly. When the regressor is also persistent --- a law that
switches on and stays on --- the observations of one unit carry much less
information than their number suggests. Classical and heteroskedasticity-robust
standard errors both assume independent observations; cluster-robust
standard errors, clustered by unit, allow any correlation within a unit and
none across units.

\begin{record}[Placebo laws]
Bertrand, Duflo and Mullainathan (2004) took state-level data on women's wages
from the Current Population Survey, about twenty years of it, and assigned
laws at random to states and dates. The laws were fictitious and could have no
effect. Estimated with state and year effects and conventional standard
errors, up to 45~\% of these placebo laws had an ``effect'' significant at
5~\%. % ledger: b3pd:bdm04.placebo
Clustering the standard errors by state, or collapsing each state's data into a
before and an after period, brought the rejection rate close to the nominal 5~\%
when the number of states was large. The paper changed the default: in a panel,
cluster at the level at which the treatment is assigned.
\end{record}

\Cref{lst:b3:panel-data:cluster} repeats the experiment on simulated data with
50 states, 20 years and errors whose persistence is 0.8.

\omcode{code/bachelor-3/02-panel-data/listings/cluster_sim.py}{1}{30}{Placebo laws on a simulated panel; \texttt{ols\_cluster} returns the three kinds of standard error.}\label{lst:b3:panel-data:cluster}

\omoutput{code/bachelor-3/02-panel-data/outputs/cluster_sim.txt}{500 placebo laws.}\label{out:b3:panel-data:cluster}

Conventional and robust standard errors reject the true zero in about 40~\% of
the samples; clustered errors in 6.6~\%. % output: cluster_sim
(The listing divides by $n - 1$ rather than by the $n - 70$ degrees of freedom left
after the state and year means are removed, which makes the conventional and robust
errors about 3.5~\% too small and flatters their rejection rates a little; the
clustered rate is barely affected.)
Clustering is not free: it needs many clusters, since its justification is
asymptotic in the number of clusters. With fewer than about thirty or forty,
clustered standard errors are themselves too small, and methods that resample
whole clusters are used instead.

\section{Dynamic panels}\label{sec:b3:panel-data:dynamic}

Many economic variables depend on their own past: this year's employment on last
year's, this year's income on last year's. The natural model is
\[
  y_{it} = \rho\, y_{i,t-1} + a_i + u_{it}, \qquad |\rho| < 1.
\]
Here strict exogeneity cannot hold: the regressor $y_{i,t-1}$ depends on
$u_{i,t-1}$, which sits inside $\bar u_i$, which sits inside the demeaned
error.

\begin{definition}[Dynamic panel bias]\label{def:b3:panel-data:nickell}
\index{dynamic panel bias}
The \emph{dynamic panel bias} (Nickell bias) is the inconsistency, as $n$ grows
with $T$ fixed, of the \omterm{def:b3:panel-data:within}{within estimator} of $\rho$ in a model with a lagged
dependent variable and unit effects.
\end{definition}

For large $T$ the bias is approximately $-(1 + \rho)/(T - 1)$ (Nickell 1981).
\admitted
The exact expression is a ratio of two sums of geometric series; its
derivation, and the approximation, are in Nickell's paper. The chart shows a
simulation.

\begin{omchart}{The within estimate of persistence in a dynamic panel with unit
effects, against the number of periods (true $\rho = 0.5$; 500 units; mean of
200 simulated panels). Dashed: the approximation $\rho - (1+\rho)/(T-1)$.}{simulated
data, the chapter's seeded script}
\begin{tikzpicture}
\begin{axis}[omtime, xlabel={Periods in the panel, $T$}, ylabel={Estimated $\rho$},
  xmin=0, xmax=31, ymin=-0.3, ymax=0.6, legend pos=south east,
  y tick label style={/pgf/number format/.cd, fixed, precision=1}, height=0.42\linewidth]
  \addplot+[mark=*, mark size=1.4pt] table[x=T, y=estimate] {figdata/bachelor-3/02-panel-data/out/nickell.csv};
  \addplot+[dashed, no marks] table[x=T, y=approx] {figdata/bachelor-3/02-panel-data/out/nickell.csv};
  \addplot[guide, thick, domain=0:31, samples=2] {0.5};
  \legend{within estimate, approximation, true value}
\end{axis}
\end{tikzpicture}
\end{omchart}

With five periods the estimate is 0.168 instead of 0.5; with thirty, 0.448.
% figdata: nickell
In a panel of countries over decades the bias is small; in a panel of firms over
five years it is most of the estimate. The cure is the tool of
\cref{ch:b3:instrumental-variables}. In first differences, $\Delta y_{i,t-1}$ is
correlated with $\Delta u_{it}$ through $u_{i,t-1}$, but $y_{i,t-2}$, dated
before that error, is not: under no serial correlation in $u$, the level two
periods back is a valid instrument for the lagged difference. Estimators built
on this idea, using every earlier level as an instrument, are standard in
dynamic panels, and so are their weak-instrument problems when $\rho$ is close
to one.

\section{Exercises}\label{sec:b3:panel-data:exercises}

\begin{exercise}[$\star$]\label{exo:b3:panel-data:1}
Read the opening chart.
\begin{enumerate}
\item State the between and within slopes in tonnes per person per log point,
and say which comparison each one makes.
\item Give two characteristics of a country that are roughly fixed over three
decades and raise both its income and its emissions.
\item The right panel's horizontal axis runs from $-0.8$ to $0.8$. Why is it so
much narrower than the left panel's range of log incomes?
\end{enumerate}
\end{exercise}

\begin{exercise}[$\star$]\label{exo:b3:panel-data:2}
Unit A has $(x, y) = (1, 2), (2, 4), (3, 5)$ in three years, and unit B has
$(4, 9), (5, 10), (6, 12)$.
\begin{enumerate}
\item Compute the pooled slope.
\item Compute the deviations from each unit's means and the within slope.
\item Draw the six points and explain the difference.
\end{enumerate}
\end{exercise}

\begin{exercise}[$\star$]\label{exo:b3:panel-data:3}
For the data of \cref{exo:b3:panel-data:2}, compute the first differences and
the first-difference slope, without a constant. Why can the slope not be
estimated with a constant in these data?
\end{exercise}

\begin{exercise}[$\star$]\label{exo:b3:panel-data:4}
The fixed-effects model assumes that the unit effect is constant over the panel.
Describe two results, from data you could compute with the opening chart's
panel, that would show it is not.
\end{exercise}

\begin{exercise}[$\star\star$]\label{exo:b3:panel-data:5}
Using the panel of \cref{lst:b3:panel-data:co2} and the helpers
\texttt{twoway\_demean} and \texttt{ols\_cluster}, estimate the within slope with
country effects only and with country and year effects, each with conventional
and country-clustered standard errors.
\begin{enumerate}
\item Report the four standard errors and the ratio of clustered to conventional.
\item By how much did the cross-country mean of log income rise between 1990
and 2022? Use it to explain why the year effects change the slope.
\end{enumerate}
\end{exercise}

\begin{exercise}[$\star\star$]\label{exo:b3:panel-data:6}
Suppose that, among countries, the richer ones have cleaner technologies for
reasons that do not change over the panel, so $\Cov(\bar x_i, a_i) < 0$ in a
regression of emissions on log income.
\begin{enumerate}
\item In which direction is the between slope biased?
\item The opening chart's between slope exceeds its within slope. Is that
consistent with this story alone? What else must be at work?
\end{enumerate}
\end{exercise}

\begin{exercise}[$\star\star$]\label{exo:b3:panel-data:7}
In \cref{exo:b3:panel-data:5}, the within slope has a conventional standard
error of 0.122 and a clustered one of 0.265.
\begin{enumerate}
\item Compute the ratio of the two variances (the design effect).
\item Dividing the 6{,}419 country-years by the design effect gives an
``effective'' number of independent observations. Compute it.
\item Give a 95~\% confidence interval for the within slope with each standard
error.
\end{enumerate}
\end{exercise}

\begin{exercise}[$\star\star$]\label{exo:b3:panel-data:8}
In \cref{out:b3:panel-data:re}, the gap between random and fixed effects is
0.003 in the first panel and 0.146 in the second.
\begin{enumerate}
\item Which estimate moves when the unit effect becomes correlated with $x$, and
why does the other not?
\item The weight $\theta$ is 0.610 in the first panel. Compute
$\sigma_u^2/(\sigma_u^2 + T\sigma_a^2)$ and compare it with the true value
when $\sigma_u^2 = \sigma_a^2 = 1$ and $T = 6$.
\end{enumerate}
\end{exercise}

\begin{exercise}[$\star\star$]\label{exo:b3:panel-data:9}
Repeat \cref{lst:b3:panel-data:co2} with the total fertility rate
(\texttt{SP.DYN.TFRT.IN}) in place of emissions, for 1990--2022.
\begin{enumerate}
\item Report the between and within slopes.
\item Unlike emissions, the two slopes are close. What does that suggest about
the fixed differences between countries, for this variable?
\end{enumerate}
\end{exercise}

\begin{exercise}[$\star\star\star$]\label{exo:b3:panel-data:10}
Prove that with $T = 2$ the \omterm{def:b3:panel-data:within}{within estimator} and the \omterm{def:b3:panel-data:first-difference}{first-difference estimator}
(without a constant) are equal.
\end{exercise}

\begin{exercise}[$\star\star\star$]\label{exo:b3:panel-data:11}
Simulate the dynamic panel of \cref{sec:b3:panel-data:dynamic} with $\rho = 0.5$
and 500 units, for $T = 5$, 10 and 20.
\begin{enumerate}
\item Report the mean within estimate over 200 panels and its bias.
\item Compare the bias with $-(1 + \rho)/(T - 1)$. For which $T$ is the
approximation good?
\end{enumerate}
\end{exercise}

\begin{exercise}[$\star\star\star$]\label{exo:b3:panel-data:12}
Strict exogeneity can fail through feedback: a shock to $y$ this year changes
$x$ next year.
\begin{enumerate}
\item Show that under feedback, $\ddot x_{it}$ is correlated with $\ddot u_{it}$
even when $x_{it}$ is uncorrelated with $u_{it}$.
\item Propose a test: add the lead $x_{i,t+1}$ to the within regression. Under
strict exogeneity, what should its coefficient be? Under feedback, what sign do
you expect if $u$ raises future $x$?
\end{enumerate}
\end{exercise}

\section{Problem: Fifteen Years of Brenland's Firms}\label{sec:b3:panel-data:problem}

\begin{problem}[{Weekend problem --- fifteen years of Brenland's firms}]\label{pb:b3:panel-data:firms}
Brenland (invented) keeps a register of 300 manufacturing firms, observed every
year from 2011 to 2025: 4{,}500 firm-years. For each firm-year it records the
logarithm of output $y_{it}$ and of capital $k_{it}$. The quality of a firm's
management, $m_i$, is not recorded; it raises output with a coefficient of 0.5
per standard deviation and is the same in every year. In 2018 a training
subsidy was given to 100 firms, chosen by a panel that judged, among other
things, how well they were run; $s_{it} = 1$ for those firms from 2018 on. The
table gives the estimates of a statistician, with standard errors computed
twice: as if all firm-years were independent, and clustered by firm. In the
register, the variance of $k_{it}$ is 0.8675 and its covariance with $m_i$ is
0.9512. % figdata: firms_facts
\begin{center}
\pgfplotstabletypeset[col sep=comma, string type,
  columns/model/.style={column name={Model}, column type=l},
  columns/coef/.style={column name={Variable}, column type=c},
  columns/estimate/.style={column name={Estimate}, column type=c},
  columns/se_iid/.style={column name={s.e. (iid)}, column type=c},
  columns/se_cluster/.style={column name={s.e. (clustered)}, column type=c},
  every head row/.style={after row=\hline}, every last row/.style={after row=\hline}]
  {figdata/bachelor-3/02-panel-data/out/firms_results.csv}
\end{center}
Here ``within'' uses firm effects, ``two-way'' firm and year effects, and
``first difference'' regresses $\Delta y$ on $\Delta k$.

\medskip\noindent\textbf{Part I --- the pooled regression.}
\begin{enumerate}
\item Interpret the pooled coefficient on $k$ in words, as an elasticity.
\item Use the omitted-variable formula and the register's variance and
covariance to predict the probability limit of the pooled slope, if the true
elasticity is 0.30.
\item Compare your prediction with the pooled estimate.
\item Why can the statistician not simply add management quality to the
regression?
\item If better-run firms used less capital, in which direction would the
pooled estimate be biased?
\item Compute the ratio of the clustered to the iid standard error of the pooled
slope. Is a precise standard error a reason to trust the estimate?
\end{enumerate}

\medskip\noindent\textbf{Part II --- within and first differences.}
\begin{enumerate}[resume]
\item Interpret the within estimate, and say which comparisons it makes.
\item The within and first-difference estimates differ by about 0.05. Give two
reasons why they may differ.
\item The two-way estimate is closer to 0.30 than the within estimate. What
does the year effect remove?
\item Give a 95~\% confidence interval for the two-way coefficient on $k$ with
the clustered standard error.
\item What part of the variation in capital identifies the two-way estimate?
\item A consultant advised some firms from 2015 and, she says, raised both their
capital and their output. Would the two-way estimate still be consistent?
Explain.
\end{enumerate}

\medskip\noindent\textbf{Part III --- the training subsidy.}
\begin{enumerate}[resume]
\item Why is the pooled estimate of the subsidy's effect likely to be biased
upward?
\item Compare the pooled and two-way estimates of the subsidy's effect.
\item Convert the two-way estimate into a percentage change in output.
\item With firm and year effects, which comparison identifies the subsidy's
effect?
\item What must be true of the treated and untreated firms' output paths for
that comparison to give the effect of the subsidy?
\item How would you use the years 2011--2017 to check that assumption?
\item Suppose the panel had picked firms whose output had just fallen. In which
direction would the two-way estimate be biased?
\end{enumerate}

\medskip\noindent\textbf{Part IV --- the standard errors.}
\begin{enumerate}[resume]
\item Compute the ratio of the clustered to the iid standard error of the
two-way subsidy estimate, and its square.
\item Give 95~\% confidence intervals for the subsidy's effect with each standard
error.
\item Why is the iid standard error too small here? Name the two features of
the data that combine.
\item If the design effect (your square from Question~20) is $1 + (T - 1)\rho$
with $T = 15$, what correlation $\rho$ does it imply?
\item How many independent firm-years is the panel worth for this coefficient?
\item State the two results of this problem, each with its clustered 95~\%
interval.
\end{enumerate}
\end{problem}
